
\chapter{Testes}
\label{chap:testes}

% Apresentar os testes realizados para verificar o funcionamento
% dos componentes e do sistema DoseCerta integrado.
% Informar o objetivo geral da etapa de testes.

\section{Objetivos e Estratégia de Testes}
\label{sec:objetivos-testes}

% Descrever o que os testes buscavam verificar.
% Explicar a estratégia adotada: testes de componentes,
% integração e sistema, conforme realizados.

\section{Ambiente e Recursos de Teste}
\label{sec:ambiente-testes}

% Informar o ambiente, equipamentos, ferramentas,
% versões de software e condições relevantes.
% Registrar limitações do ambiente, se houver.

\section{Testes do Projeto Mecânico}
\label{sec:testes-mecanicos}

% Descrever os testes aplicados ao mecanismo:
% montagem, movimentação, alinhamento, acomodação das doses
% ou outros aspectos que tenham sido efetivamente verificados.
%
% Para cada teste, registrar:
% objetivo, procedimento, critério de aceitação e resultado.

\section{Testes Eletrônicos e do Sistema Embarcado}
\label{sec:testes-eletronicos-embarcados}

% Descrever os testes de sensores, atuadores, alimentação,
% firmware e controle do mecanismo.
% Registrar instrumentos e condições de medição, quando aplicável.

\section{Testes do Backend e do Aplicativo}
\label{sec:testes-backend-aplicativo}

% Descrever testes de funcionalidades, validação de dados,
% comunicação entre módulos e tratamento de erros.
% Incluir testes automatizados ou manuais, conforme realizados.

\section{Testes de Visão Computacional e do Agente de IA}
\label{sec:testes-visao-ia}

% Descrever os testes realizados para avaliar os componentes
% de visão computacional e IA, caso tenham sido implementados.
% Informar entradas, saídas esperadas, critérios e resultados.
% Não declarar métricas sem dados coletados.

\section{Testes de Integração}
\label{sec:testes-integracao}

% Verificar o funcionamento conjunto dos subsistemas.
% Descrever os cenários de integração testados,
% as condições de execução e os problemas encontrados.

\section{Testes de Usabilidade}
\label{sec:testes-usabilidade}

% Se houve avaliação com usuários, descrever participantes
% de forma apropriada, tarefas, procedimento e instrumento
% de coleta. Apresentar resultados sem extrapolar a amostra.
% Se não houve avaliação com usuários, registrar essa limitação.

\section{Registro e Análise dos Resultados dos Testes}
\label{sec:registro-resultados-testes}

% Organizar os resultados obtidos e as evidências:
% tabelas, fotografias, capturas de tela, logs ou medições.
% Explicar falhas, ajustes realizados e retestes.

% Exemplo de tabela de registro:
% \begin{table}[htbp]
%     \centering
%     \caption{Registro dos testes realizados}
%     \label{tab:registro-testes}
%     \begin{tabular}{|l|p{4cm}|p{3cm}|p{3cm}|}
%         \hline
%         Identificação & Procedimento & Critério de aceitação & Resultado \\
%         \hline
%         % T01 & Descrever o teste & Critério definido & A preencher \\
%         \hline
%     \end{tabular}
% \end{table}

\section{Síntese dos Testes}
\label{sec:sintese-testes}

% Resumir quais funções foram verificadas,
% quais atenderam aos critérios e quais permaneceram
% pendentes ou apresentaram limitações.
